% Paper 2, §11 — RL on verifiable rewards: GRPO and DAPO
% Anchors: kb/notes/post-training/rlvr-and-grpo.md
%          kb/excerpts/{shao2024,deepseek-r1,dapo2025,ouyang2022-instructgpt}.md
% Cross-paper: anchor for Paper 3 §7 (the reasoning triangle's RL leg).

\section{RL on verifiable rewards: GRPO and DAPO}
\label{sec:rlvr-grpo}
\label{sec:rlvr-frontier}

The post-training pipeline of \cref{sec:rlhf} and \cref{sec:dpo} uses a
\emph{learned} reward $r_\phi(x, y)$ — a Bradley--Terry head fit to
human preference pairs, or its \glsterm{dpo}{DPO} closed-form analog — as the gradient
signal for the policy. The \glsterm{rlvr-3}{RLVR}-era reformulation replaces that learned
reward with a deterministic verifier: $R(x, y) = \mathbb{1}[\,\mathrm{verify}(x, y) = \mathsf{true}\,]$,
the indicator that the model's output passes a sympy equivalence check,
a unit-test suite, or a regex format check
\citep[\S 2.2]{deepseek-r1}.\footnote{KB excerpt:
\texttt{kb/excerpts/deepseek-r1.md\#sec-2-2-reward}.} The thesis of this
section is that two coupled changes — a programmatic verifier in place
of a learned \glsterm{reward-model}{reward model}, and a group-mean baseline in place of a
co-sized \glsterm{value}{value} model \citep[\S 4.1]{shao2024}\footnote{KB excerpt:
\texttt{kb/excerpts/shao2024.md\#sec-4-1}.} — together produce the
post-training recipe behind \citet{deepseek-r1}, the Qwen 2.5 / 3
reasoning lines \citep{qwen3}, T\"ulu 3 \citep{tulu3}, and the
open-weight RL-for-reasoning ecosystem of 2024--2026. The headline
empirical fact is also the cleanest single number in the entire training
pipeline: \citet[\S 2.3]{deepseek-r1}\footnote{KB excerpt:
\texttt{kb/excerpts/deepseek-r1.md\#sec-2-3-aime}.} report
DeepSeek-\glsterm{r1-zero}{R1-Zero}'s average pass@1 on AIME 2024 climbing from
$15.6\%$ to $77.9\%$ over the course of one continuous \glsterm{grpo-2}{GRPO} run, and to
$86.7\%$ with \glsterm{self-consistency-2}{self-consistency} over $64$ samples — surpassing the
average human competitor. This section is also the cross-paper anchor
that \hyperref[P3-sec:rlvr-grpo-dapo]{Paper~3 §7} cites for the \emph{training} leg of its reasoning
triangle (compute / search / verification); the \emph{inference} side of
the same machinery is treated there.

\subsection{What ``verifiable'' means}
\label{sec:rlvr-verifiable}

Define the verifiable-reward setting: for prompt $x$ and sampled output
$y$, the reward is
\begin{equation}
R_{\mathrm{verifier}}(x, y) \;=\; \mathbb{1}\big[\,\mathrm{verify}(x, y) = \mathsf{true}\,\big] \;\in\; \{0, 1\},
\label{eq:verifier-reward}
\end{equation}
where $\mathrm{verify}$ is a deterministic program supplied alongside
the prompt. The instances that matter at scale
\citep[\S 2.2]{deepseek-r1}:

\begin{itemize}
\item \emph{Math.} $\mathrm{verify}$ is exact match against the canonical
answer after sympy-style normalization (rationalization,
trig-identity reduction, whitespace canonicalization). Competition-style
benchmarks (AIME, MATH, AMC) have a small set of canonical answers per
problem, almost always integers or short symbolic expressions, so the
exact-match relation is well-defined.
\item \emph{Code.} $\mathrm{verify}$ is ``compiles, then passes the
attached test cases.'' Per-problem unit-test suites are the gold reward;
competitive programming benchmarks (LiveCodeBench, Codeforces) ship
them.
\item \emph{Format.} $\mathrm{verify}$ is a parser or regex (e.g., the
trajectory contains a well-formed
\texttt{<think>...</think>} block before its answer) — a small constant
add-on reward that keeps the reasoning channel alive during early
training \citep[\S 2.2]{deepseek-r1}.
\end{itemize}

The reward is \emph{sparse} (one bit per trajectory) and \emph{terminal}
(only the final answer is graded; intermediate steps are not, in the
\glsterm{rlvr-3}{RLVR} setting; cf.\ \cref{sec:rlvr-frontier} on process supervision). The
contrast with \cref{sec:rlhf}'s \glsterm{rlhf}{RLHF} setting is structural: there,
$r_\phi$ is a Bradley--Terry head trained on $\binom{K}{2}$ human
comparisons per prompt, and the policy can game artifacts of $r_\phi$
that do not correspond to true quality (verbosity, \glsterm{sycophancy}{sycophancy}, surface
formatting). With \cref{eq:verifier-reward}, the only way to score $1$
is to actually solve the problem, which sidesteps two failure modes of
learned reward models:

\begin{enumerate}
\item \emph{\glsterm{reward-hacking}{Reward hacking}.} The policy exploits an artifact of
$r_\phi$. Programmatic verifiers admit this only when the verifier
itself has a bug (e.g., a normalization rule that accepts
$0 = 0\cdot x$ for any $x$); see \cref{sec:rlvr-frontier}.
\item \emph{\glsterm{reward-overoptimization}{Reward overoptimization}.} The proxy--gold gap grows with
$\KL(\pi_\theta \| \pi_{\mathrm{ref}})$ in \glsterm{rlhf}{RLHF} (Gao et al.\ 2023,
``Scaling laws for \glsterm{reward-model}{reward model} overoptimization'').
\deepencite{gao2023-overopt — bibkey not yet in references.bib} Verifiers do not have this gap — they
\emph{are} the gold reward, modulo dataset-side issues like wrong gold
answers.
\end{enumerate}

\citet[\S 2.2]{deepseek-r1} state this stance explicitly: ``we abstain
from applying neural reward models — whether outcome-based or
process-based — to reasoning tasks. This decision is predicated on our
observation that neural reward models are susceptible to \glsterm{reward-hacking}{reward hacking}
during large-scale reinforcement learning.''\footnote{KB excerpt:
\texttt{kb/excerpts/deepseek-r1.md\#sec-2-2-reward}.}

\subsection{GRPO: critic-free PPO with a group-mean baseline}
\label{sec:rlvr-grpo-objective}

\citet{shao2024} introduce \emph{\glsterm{grpo}{Group Relative Policy Optimization}}
(\glsterm{grpo-2}{GRPO}) in the DeepSeekMath paper as a critic-free PPO variant. The
algorithm samples a group of $G$ trajectories per prompt and uses the
group's empirical reward distribution as the variance-reducing
baseline, eliminating the \glsterm{value}{value} network $V_\phi$ entirely
\citep[\S 4.1, motivation]{shao2024}.\footnote{KB excerpt:
\texttt{kb/excerpts/shao2024.md\#sec-4-1-motivation}.} The objective,
restated and verified in \citet[\S 2.1, Eq.\ 1]{deepseek-r1}\footnote{KB
excerpts: \texttt{kb/excerpts/shao2024.md\#sec-4-1} and
\texttt{kb/excerpts/deepseek-r1.md\#sec-2-1-grpo}.}:
\begin{equation}
\begin{aligned}
\mathcal{J}_{\mathrm{GRPO}}(\theta) \;=\; & \E_{q \sim P(Q),\, \{o_i\}_{i=1}^G \sim \pi_{\theta_{\mathrm{old}}}(\cdot \mid q)} \bigg[ \\
& \quad \frac{1}{G} \sum_{i=1}^{G} \frac{1}{|o_i|} \sum_{t=1}^{|o_i|}
\min\!\Big( r_{i,t}(\theta)\, \hat{A}_{i,t},\;
\mathrm{clip}\!\big(r_{i,t}(\theta), 1 - \epsilon, 1 + \epsilon\big)\, \hat{A}_{i,t} \Big) \\
& \quad -\, \beta \, \KL\!\big[\pi_\theta \,\|\, \pi_{\mathrm{ref}}\big] \bigg],
\end{aligned}
\label{eq:grpo}
\end{equation}
with importance ratio
$r_{i,t}(\theta) = \pi_\theta(o_{i,t} \mid q, o_{i,<t}) \,/\, \pi_{\theta_{\mathrm{old}}}(o_{i,t} \mid q, o_{i,<t})$
and \emph{\glsterm{group-normalized-advantage}{group-normalized advantage}}
\begin{equation}
\hat{A}_{i,t} \;=\; \tilde{A}_i \;=\; \frac{R_i \,-\, \mathrm{mean}\!\big(\{R_1, \ldots, R_G\}\big)}{\mathrm{std}\!\big(\{R_1, \ldots, R_G\}\big)},
\label{eq:grpo-advantage}
\end{equation}
constant in $t$ across the trajectory \citep[\S 4.1, Eq.\ verified
against R1 \S 2.1 Eq.\ 3]{shao2024}. The KL penalty is computed with
Schulman's unbiased estimator,
$\KL = \tfrac{\pi_{\mathrm{ref}}}{\pi_\theta} - \log\tfrac{\pi_{\mathrm{ref}}}{\pi_\theta} - 1$,
evaluated per token \citep[\S 4.1]{shao2024}.

\paragraph{Symbol table.} $q$ is the prompt; $G$ is the group size,
typically $8$--$64$; $o_i$ is the $i$-th sampled trajectory
(\glsterm{chain-of-thought}{chain-of-thought} + answer); $R_i \in \{0, 1\}$ is the verifier reward;
$\pi_{\mathrm{ref}}$ is a frozen reference policy (typically the SFT
cold-start, cf.\ \cref{sec:sft}); $\beta$ is the KL coefficient,
typically $0.001$--$0.04$ — \citet[\S 2.1]{deepseek-r1} report
$\beta = 0.001$ for \glsterm{r1-zero}{R1-Zero};\footnote{KB excerpt:
\texttt{kb/excerpts/deepseek-r1.md\#sec-2-1-hp}.} $\epsilon$ is the PPO
clip range, typically $0.2$.

\paragraph{Memory profile (the load-bearing reason GRPO matters at LLM
scale).} PPO with GAE \citep[\S 3.5]{ouyang2022-instructgpt} requires a
\glsterm{value}{value} network $V_\phi(s_t)$ co-sized with the policy: at \glsterm{llm}{LLM} scale this
\emph{doubles} the active parameter count and the activation memory of
training. \glsterm{grpo-2}{GRPO}'s group-mean baseline replaces $V_\phi$ entirely; the
trade is $G \times$ rollout count per prompt against $V_\phi$'s
parameter and activation footprint, and at sparse reward (where $V_\phi$
struggles to learn anyway) the trade is favorable
\citep[\S 4.1]{shao2024}.\footnote{KB excerpt:
\texttt{kb/excerpts/shao2024.md\#sec-4-1-motivation}.} The four
co-resident model copies of \cref{sec:rlhf}'s PPO infrastructure
(policy, reference, reward, \glsterm{value}{value}) reduce to two (policy, reference);
the verifier is a CPU subprocess, not a model.

\begin{intuition}
\glsterm{grpo-2}{GRPO} is the median-of-rollouts as a policy gradient. Sample $G$
trajectories from the current policy, score each, and ask: which were
\emph{better than typical for this prompt}? Up-weight the
better-than-mean log-probs; down-weight the worse-than-mean ones. The
sign of $\tilde{A}_i$ in \cref{eq:grpo-advantage} answers the question;
its magnitude (in standard-deviation units of the group's reward
distribution) controls the step size. Returning to canonical form:
\cref{eq:grpo-advantage} is the standardization of $R_i$ within its
group, and \cref{eq:grpo} is the PPO-clipped surrogate with that
standardized advantage broadcast over the trajectory's tokens.
\end{intuition}

\paragraph{Per-step procedure.} For each batch of $B$ prompts: (1)
sample $G$ trajectories per prompt from $\pi_{\theta_{\mathrm{old}}}$ at
temperature $T \in [0.7, 1.0]$, capped at $L_{\max}$ tokens (DeepSeek-R1
uses $L_{\max} = 32{,}768$ before step $8.2$k and $65{,}536$
afterward \citep[\S 2.1]{deepseek-r1}); (2) score with the verifier;
(3) compute $\tilde{A}_i$ via \cref{eq:grpo-advantage} and broadcast
token-wise; (4) take $K \in \{1, 4\}$ PPO inner-loop steps on the
buffered $(B \cdot G)$ trajectories; (5) refresh
$\theta_{\mathrm{old}} \leftarrow \theta$.

\subsection{R1 evidence: 15.6\% to 77.9\% on AIME 2024}
\label{sec:rlvr-r1}

\citet{deepseek-r1} apply \glsterm{grpo-2}{GRPO} with verifiable rewards to the
\glsterm{precision-pinning}{DeepSeek-V3} base model and report the strongest published evidence to
date that pure RL on a deterministic verifier produces frontier-grade
reasoning capability. Two models are trained:

\begin{itemize}
\item \emph{DeepSeek-\glsterm{r1-zero}{R1-Zero}.} \glsterm{grpo-2}{GRPO} applied directly to the base, no SFT
cold-start. Reward is rule-based: accuracy (binary correctness from a
math/code verifier) plus a small \glsterm{format-reward}{format reward} for the
\texttt{<think>...</think><answer>...</answer>} envelope
\citep[\S 2.2]{deepseek-r1}.\footnote{KB excerpt:
\texttt{kb/excerpts/deepseek-r1.md\#sec-2-2-reward}.}
\item \emph{DeepSeek-R1.} A multi-stage pipeline: cold-start SFT on a
small set of long-\glsterm{chain-of-thought}{CoT} exemplars, then \glsterm{grpo-2}{GRPO} with the same verifiable
rewards (plus a language-consistency reward), then a second SFT on
rejection-sampled reasoning traces, then a final RL alignment stage for
helpfulness and harmlessness \citep[\S 3]{deepseek-r1}.\footnote{KB
excerpt: \texttt{kb/excerpts/deepseek-r1.md\#sec-3-pipeline}.}
\end{itemize}

The training-curve headline \citep[\S 2.3]{deepseek-r1}\footnote{KB
excerpt: \texttt{kb/excerpts/deepseek-r1.md\#sec-2-3-aime}.}: ``Figure
1(a) depicts the performance trajectory of DeepSeek-\glsterm{r1-zero}{R1-Zero} on the AIME
2024 benchmark throughout the RL training process, where the average
pass@1 score on AIME 2024 shows a significant increase, jumping from an
initial $15.6\%$ to $77.9\%$. In addition, by leveraging the
\glsterm{self-consistency-2}{self-consistency} decoding, the model's performance can be further
improved, achieving an accuracy of $86.7\%$. This performance
significantly surpasses the average performance across all human
competitors.'' The full run is $10{,}400$ steps,
$\sim\!1.6$ epochs over the training distribution
\citep[\S 2.1]{deepseek-r1}.

Two structural facts of the trajectory are load-bearing for the
section's thesis:

\paragraph{Response length grows during training.} The output trace
length grows from a few hundred tokens early in training to roughly
$10{,}000$ tokens at convergence, accompanied by the emergence of
self-reflective patterns
\citep[\S 2.3]{deepseek-r1}.\footnote{KB excerpt:
\texttt{kb/excerpts/deepseek-r1.md\#sec-2-3-aha}.} The verifier reward
contains no length term — the policy discovers that long traces win on
math and code, on its own. Cross-paper to Paper~3 §4 for the
inference-time-compute interpretation of this growth.

\paragraph{The ``aha moment.''} \citet[\S 2.3]{deepseek-r1} report a
discrete jump in the model's use of the word ``\texttt{wait}'' during
reflection, characterized as ``an `\glsterm{aha-moment}{aha moment}', characterized by a
sudden increase in the use of the word `wait' during reflections.''
\footnote{KB excerpt:
\texttt{kb/excerpts/deepseek-r1.md\#sec-2-3-aha}.} The \glsterm{chain-of-thought}{chain-of-thought}
structure is being actively reorganized during RL, not just used.

\paragraph{R1 vs.\ R1-Zero — what the cold-start adds.} \glsterm{r1-zero}{R1-Zero}'s traces
are difficult to read and language-mix between English and Chinese
because \glsterm{precision-pinning}{DeepSeek-V3}-Base is multilingual
\citep[\S 3]{deepseek-r1}.\footnote{KB excerpt:
\texttt{kb/excerpts/deepseek-r1.md\#sec-3-pipeline}.} R1's cold-start
SFT adds: (a) legibility (single-language English traces with consistent
formatting); (b) the \texttt{<think>...</think>} envelope as a learned
default. R1's reasoning capability matches \glsterm{r1-zero}{R1-Zero}'s; the cold-start is
read by the authors as primarily stylistic.

\paragraph{Distillation transfer (cross-link to \cref{sec:adaptation}
and Paper~3 §8).} R1 generates $800$K reasoning traces (600K reasoning
+ 200K non-reasoning), and \emph{pure SFT} on this dataset transfers
reasoning capability into Qwen-class small models — distilled-32B
matches or exceeds OpenAI o1-mini on multiple reasoning benchmarks
\citep[\S 4]{deepseek-r1}.\footnote{KB excerpt:
\texttt{kb/excerpts/deepseek-r1.md\#sec-4}.} The capability is
expressible as an SFT dataset; small-model trainers do not need to run
RL.

\subsection{DAPO: four targeted refinements}
\label{sec:rlvr-dapo}

\citet{dapo2025} (``Decoupled clip and dynamic s\textbf{A}mpling
\textbf{P}olicy \textbf{O}ptimization,'' ByteDance Seed, March 2025)
introduce four targeted modifications to \glsterm{grpo-2}{GRPO} that together cut
training-step count to AIME-2024 score 50 by roughly half on
Qwen2.5-32B \citep[\S 4]{dapo2025}.\footnote{KB excerpt:
\texttt{kb/excerpts/dapo2025.md\#sec-4-headline}.} Each modification
addresses a measured failure mode of the \glsterm{grpo-2}{GRPO} recipe.

\paragraph{(i) Clip-Higher (Eq.\ 10).} \glsterm{grpo-2}{GRPO}'s symmetric clip
$[1 - \epsilon, 1 + \epsilon]$ caps probability \emph{increases} on
exploration tokens at the same rate as it caps decreases on
already-favored tokens. \glsterm{dapo-2}{DAPO} decouples the bound:
\begin{equation}
\mathcal{J}_{\mathrm{DAPO}}(\theta) \;=\; \E\!\left[\frac{1}{G} \sum_i \min\!\Big( r_{i,t}(\theta)\, \hat{A}_i,\; \mathrm{clip}\!\big(r_{i,t}(\theta), 1 - \epsilon_{\mathrm{low}}, 1 + \epsilon_{\mathrm{high}}\big)\, \hat{A}_i \Big) \right],
\label{eq:dapo-clip}
\end{equation}
with $\epsilon_{\mathrm{low}} = 0.2$ kept at the PPO default and
$\epsilon_{\mathrm{high}}$ enlarged (e.g., $0.28$)
\citep[Eq.\ 10]{dapo2025}.\footnote{KB excerpt:
\texttt{kb/excerpts/dapo2025.md\#sec-3-clip-higher}.} The asymmetry
permits larger upward updates on under-probable correct tokens — i.e.,
prevents collapse onto already-favored trajectories.

\paragraph{(ii) Dynamic Sampling (Eq.\ 11).} Whenever all $G$
trajectories in a group share the same reward (all 0 or all 1), the
\glsterm{group-normalized-advantage}{group-normalized advantage} \cref{eq:grpo-advantage} is identically zero
and the sample contributes no gradient. \glsterm{dapo-2}{DAPO} enforces
\begin{equation}
0 \;<\; \big| \{\, o_i : \mathrm{is\_equivalent}(a, o_i) \,\}\big| \;<\; G
\label{eq:dapo-dynsample}
\end{equation}
\citep[Eq.\ 11]{dapo2025}, discarding and re-sampling groups that fail
the constraint.\footnote{KB excerpt:
\texttt{kb/excerpts/dapo2025.md\#sec-3-dynamic-sampling}.} Compute is
concentrated on informative gradients.

\paragraph{(iii) Token-Level Loss (Eq.\ 12).} \glsterm{grpo-2}{GRPO} normalizes by
$1/|o_i|$ inside the group sum and then by $1/G$ outside, which means
longer trajectories contribute proportionally less per token. \glsterm{dapo-2}{DAPO}
replaces this with a token-aggregated normalization:
\begin{equation}
\mathcal{L}_{\mathrm{DAPO}}(\theta) \;=\; -\,\frac{1}{\sum_{i=1}^{G} |o_i|} \sum_{i=1}^{G} \sum_{t=1}^{|o_i|} L_{i,t}(\theta).
\label{eq:dapo-tokenloss}
\end{equation}
Every token weighs equally regardless of trajectory length
\citep[Eq.\ 12]{dapo2025}.\footnote{KB excerpt:
\texttt{kb/excerpts/dapo2025.md\#sec-3-token-loss}.} The change avoids
the systematic underweighting of long correct CoTs — which is precisely
the regime \cref{sec:rlvr-r1}'s response-length growth produces.

\paragraph{(iv) Overlong Reward Shaping (Eq.\ 13).} Trajectories that
hit $L_{\max}$ are typically penalized as failures, but the failure
signal is uninformative — the model just ran out of budget. \glsterm{dapo-2}{DAPO}
replaces the hard zero with a length-based soft penalty:
\begin{equation}
R_{\mathrm{length}}(o_i) \;=\; \begin{cases}
0 & |o_i| \le L_{\mathrm{accept}} \\
\text{linear interp.} & L_{\mathrm{accept}} < |o_i| < L_{\max} \\
-1 & |o_i| \ge L_{\max}
\end{cases}
\label{eq:dapo-overlong}
\end{equation}
\citep[Eq.\ 13]{dapo2025}.\footnote{KB excerpt:
\texttt{kb/excerpts/dapo2025.md\#sec-3-overlong}.} Variance from
length-truncated samples falls; the genuine signal that an overlong
reasoning trajectory is bad is preserved.

\subsection{Why verifiable rewards changed the field}
\label{sec:rlvr-why}

The pre-\glsterm{rlvr-3}{RLVR} post-training stack (\cref{sec:rlhf}, \cref{sec:dpo}) was
constrained by the \emph{reward-hacking ceiling}: as
$\KL(\pi_\theta \| \pi_{\mathrm{ref}})$ grew, the proxy--gold gap of
$r_\phi$ widened (Gao et al.\ 2023, scaling laws for reward-model
overoptimization), and the policy began to
exploit artifacts of $r_\phi$ rather than improve along the gold-reward
direction. Mitigations (RM ensembles, \glsterm{warm}{WARM}/\glsterm{warp}{WARP}, KL-coefficient
scheduling, early stopping) bought margin but did not abolish the
ceiling. Three structural changes follow from \cref{eq:verifier-reward}:

\begin{enumerate}
\item \emph{The proxy--gold gap closes by construction.} The verifier
\emph{is} the gold reward in the verifier-supported domains, modulo
dataset-side issues. Reward-overoptimization curves of the
$r_\phi$-trained-on-pairs kind do not exist in this regime.
\item \emph{Human preference data is not the bottleneck.} A canonical
preference dataset costs $\$10$--$\$30$ per pair at frontier annotator
quality \citep[\S 3.4]{ouyang2022-instructgpt}; a verifier costs the
engineering time to write the parser plus the gold-answer set. The
training-data axis decouples from the labeling-budget axis.
\item \emph{Capability transfers through SFT distillation.}
\cref{sec:rlvr-r1}'s 800K-trace distill-into-Qwen result is the
2025 demonstration that capability gained under verifier-driven RL is
expressible as an SFT dataset, and the small-model trainer does not
need to run RL at all
\citep[\S 4]{deepseek-r1}.\footnote{KB excerpt:
\texttt{kb/excerpts/deepseek-r1.md\#sec-4}.} The post-training-stack
output of one lab becomes the SFT input of another.
\end{enumerate}

\begin{intuition}
Verifiable rewards convert post-training from a preference-elicitation
problem to a search problem. The verifier is a pass/fail oracle, not a
teacher: \glsterm{rlvr-3}{RLVR} does not tell the policy \emph{how} to be better, only
\emph{whether} it is. The optimization works because the policy already
produces the correct answer some fraction $p_0 > 0$ of the time at the
start of training; \glsterm{rlvr-3}{RLVR} amplifies the trajectories that lead to correct
answers and suppresses the rest. Returning to canonical form: the group
baseline in \cref{eq:grpo-advantage} is exactly the empirical estimate
of $p_0(q)$ for prompt $q$, and \cref{eq:grpo}'s gradient becomes a
trajectory-level importance-weighted update on the
better-than-empirical-mean rollouts.
\end{intuition}

\subsection{Cross-paper anchor: this is also Paper~3's mechanism}
\label{sec:rlvr-paper3-link}

This section is also Paper~2's anchor for Paper~3's reasoning triangle.
The triangle (Paper~3 §1, §7) has three legs: inference-time compute
scaling (\glsterm{chain-of-thought}{chain-of-thought}, \glsterm{self-consistency-2}{self-consistency}, sampling-based test-time
compute), inference-time verifier signals (process reward models, beam
/ MCTS search, generative PRMs), and \emph{RL on verifiable rewards} —
the leg developed here. Paper~3 §7 cites
\cref{eq:verifier-reward}, \cref{eq:grpo}, \cref{eq:grpo-advantage}, and
the four \glsterm{dapo-2}{DAPO} equations as the canonical statement of the \glsterm{rlvr-3}{RLVR} leg, and
develops the inference-time interactions (proposer/verifier, MCTS over
reasoning steps, ReST-MCTS* AlphaZero-loop) that this section does not
cover. The R1 trajectory of \cref{sec:rlvr-r1} is the empirical proof
point that all three legs co-design: the same \glsterm{rlvr-3}{RLVR} run produces both
the response-length growth (inference-time compute) and the
self-reflection patterns (which interact with verifier signals at
inference). Treating \glsterm{rlvr-3}{RLVR} as a training-only mechanism — as this
section deliberately does — understates the gains and, more
importantly, the open questions; the latter are concentrated in
Paper~3.

\begin{contradiction}
\textbf{Does \glsterm{rlvr-3}{RLVR} create new reasoning capability?} The framing in
\citet{deepseek-r1} (``incentivizing reasoning capability,'' the abstract
phrasing) and the parallel framings in \citet{tulu3} and \citet{qwen3}
treat \glsterm{rlvr-3}{RLVR} as eliciting reasoning behaviors that the base model could
not produce. \citet{rlvr-limits-2025} (Yue et al.\ NeurIPS 2025 oral)
disagree: by carefully comparing pass@1 vs.\ pass@$K$ of base-model
vs.\ \glsterm{rlvr-3}{RLVR}'d model at matched sampling budget, they argue \glsterm{rlvr-3}{RLVR} improves
pass@1 (sampling efficiency) on problems the base model can already
solve at higher $K$, but does \emph{not} expand the set of problems
solvable at any $K$. The \glsterm{capability-ceiling}{capability ceiling} is set by pre-training;
\glsterm{rlvr-3}{RLVR} is sampling-efficient sharpening, not capability expansion. As of
2026-Q2 the open-source consensus leans toward Yue with caveats: the
result is benchmark-, temperature-, and sample-budget-conditional, and
generalization to non-math (e.g., agentic, long-horizon) \glsterm{rlvr-3}{RLVR} is
single-source. Cross-paper to Paper~3 §14 for the resolution attempt.
\end{contradiction}

\begin{contradiction}
\textbf{\glsterm{cold-start-sft-2}{Cold-start SFT} necessity.} \glsterm{r1-zero}{R1-Zero} (no cold-start) reaches AIME
$77.9\%$ pass@1 with unreadable traces; R1 (with cold-start) is cleaner
\citep[\S 3]{deepseek-r1}. Whether the cold-start is
capability-load-bearing or only stylistic is open: anecdotal evidence
from 2025--2026 reproductions (per the topics-survey sweep recorded in
\texttt{kb/notes/post-training/\glsterm{capability-ceiling}{rlvr}-and-grpo.md\#5-4}) shows removing
cold-start sometimes destabilizes RL, sometimes works fine; the
determining factor is unclear.
\end{contradiction}

\begin{contradiction}
\textbf{\glsterm{vapo}{VAPO} vs.\ \glsterm{dapo-2}{DAPO}.} \glsterm{vapo}{VAPO} (\glsterm{vapo}{Value-Augmented PPO}, ByteDance, April
2025, arXiv:2504.05118) re-introduces a learned \glsterm{value}{value}
function under the \glsterm{grpo-2}{GRPO}/\glsterm{rlvr-3}{RLVR} objective, shared-backbone with the
policy and \glsterm{value}{value}-loss-clipped, and reports outperforming \glsterm{dapo-2}{DAPO} on
long-\glsterm{chain-of-thought}{CoT} tasks in fewer steps — suggesting the \glsterm{grpo-2}{GRPO}-vs-PPO advantage is
not ``\glsterm{value}{value} models are useless'' but ``na\"ive \glsterm{value}{value} models are
expensive-and-useless; carefully-trained ones are useful.'' Whether
\glsterm{vapo}{VAPO}'s gain over \glsterm{dapo-2}{DAPO} replicates outside the original recipe is open as
of 2026-Q2; only the original paper has reported the comparison.
\deepencite{vapo2025 — bibkey not yet in references.bib; VAPO §3 replication evidence}
\end{contradiction}

\subsection{Forward link}
\label{sec:rlvr-forward}

\cref{sec:rlvr-grpo-objective}--\cref{sec:rlvr-dapo} concludes the
preference-and-verifier RL stage of the post-training pipeline.
The recipe is now: SFT cold-start (\cref{sec:sft})
$\to$ optionally a preference-RL stage (\cref{sec:rlhf} or
\cref{sec:dpo} for instruction-following and alignment)
$\to$ \glsterm{rlvr-3}{RLVR} with \glsterm{grpo-2}{GRPO} (or \glsterm{dapo-2}{DAPO}) for reasoning capability on
verifier-supported domains $\to$ \glsterm{rejection-sampling-sft}{rejection-sampling SFT} to consolidate.
What remains is the weight-space tail of the pipeline: how to take the
post-trained checkpoint and adapt it cheaply (\glsterm{lora}{LoRA}, \glsterm{qlora}{QLoRA}), or compose
several post-trained checkpoints into a multi-task model (Soups, Task
Vectors, TIES). \cref{sec:adaptation} walks the adaptation-and-merging
machinery that closes the six-stage pipeline of this paper.
